\subsection{Severi Tommaso} \label{sec:implementation-severi}
Mi sono occupato della progettazione e implementazione di componenti appartenenti a ogni
dominio del progetto, nello specifico mi sono occupato di:
\begin{itemize}
    \item Modellazione di Point2D
    \item Modellazione dei Pawns e dunque delle Pocket e dei Collectables
    \item Implementazione del PartyController e dei relativi managers per gestire il tiro dei dadi
          iniziale e il susseguirsi dei turni
    \item Implementazione delle classi relative al Pub-Sub
    \item Implementazione della pagina finale del Party
    \item Ogni aspetto implementativo riguardante il minigioco Labirinto
    \item Scrittura logica Scala per la generazione e risoluzione del labirinto
\end{itemize}

Nelle sezioni seguenti ci sono quegli aspetti che ritengo più significativi e che espandono o si
discostano maggiormente dal design iniziale, sfruttando le potenzialità di Scala.

\subsubsection{Point2D Generico}
Sebbene non sia un aspetto principale è sicuramente un interessante esempio di come si possano
sfruttare le potenzialità del linguaggio. Si serve di due costrutti principali in modo da rimanere
il quanto più generico possibile rispetto alla rappresentazione interna:
\begin{lstlisting}[language=Scala]
trait Point2D[N](using num: Numeric[N]):
    def x: N
    def y: N
    def +(other: Point2D[N]): Point2D[N]
    def -(other: Point2D[N]): Point2D[N]
    def isZero: Boolean
    def distanceFrom(other: Point2D[N]): Double
\end{lstlisting}
Grazie all'uso del tipo generico e del type class Numeric, è possibile
utilizzare qualsiasi tipo numerico per rappresentare le coordinate, non dovendo riscrivere 
il codice per ogni tipo. Inoltre, grazie all'uso di using, si possono utilizzare le operazioni 
aritmetiche sui numeri senza dover specificare il tipo numerico ogni volta.
Sfruttando, infine, la conversione implicita, un Point2D può essere facilmente convertito in una 
tupla e viceversa senza usare cast espliciti:
\begin{lstlisting}[language=Scala]
given [N](using num: Numeric[N]): Conversion[(N, N), Point2D[N]] with
    def apply(t: (N, N)): Point2D[N] =
      Point2D(t._1, t._2)

given [N](using Numeric[N]): Conversion[Point2D[N], (N, N)] with
  def apply(p: Point2D[N]): (N, N) =
    (p.x, p.y)
\end{lstlisting}
Questo ha permesso poi di poter creare, a seconda del gioco, il proprio tipo di Point2D. Anche
grazie ai tipi opachi di Scala, usando i quali non è necessario esporre il tipo numerico scelto, 
come ad esempio nel caso di MazePosition:
\begin{lstlisting}[language=Scala]
object MazePosition:
  opaque type MazePosition = Point2D[Int]
  def apply(x: Int, y: Int): MazePosition = Point2D(x, y)
  extension (bp: MazePosition)
    def x: Int = bp.x
    def y: Int = bp.y
    def toPoint2D: Point2D[Int] = bp
    def +(dir: Direction): MazePosition =
      val (dx, dy) = dir.offset
      MazePosition(bp.x + dx, bp.y + dy)
\end{lstlisting}
Grazie poi agli extension methods, è possibile decidere quali funzionalità di Point2D mostrare.

\subsubsection{Pocket e Collectables Funzionali}
Sebbene non fosse pensato nel design iniziale, ho decido di porre una differenza tra i Collectbles:
quelli definiti come oggetto univoco e il gruppo a cui appartengono. 
\begin{lstlisting}[language=Scala]
enum Collectable:
  case Monad()
  case Rung(monadsNeeded: Int)

enum CollectableType:
  case MonadType, RungType  
\end{lstlisting}
Questa distinzione nasce dalla natura delle Pocket, in quanto non pongono nessuna differenza
negli oggetti collezionati, la loro univocità è garantita solo fino a che si trovano sul tabellone.
Per questo motivo ho deciso di implementare una Pocket che potesse effettuare operazioni sui gruppi
di oggetti collezionabili, come ad esempio rimuovere un certo numero di oggetti di un certo tipo
o aggiungerne un certo numero. Questa funzionalità non è stata sfruttata a pieno nel progetto attuale,
se non a livello di View, ma ho ritenuto potesse essere utile per sviluppi futuri. Per esempio,
se si volesse implementare l'idea di poter perdere uno, o più, oggetti dopo averli raccolti.
\newpage
\begin{lstlisting}[language=Scala]
trait Pocket:
  def getAll: Seq[Collectable]
  def getByType(itemType: CollectableType): Seq[Collectable]
  def add(item: Collectable): Pocket
  def remove(item: Collectable): Pocket
  def addMultiple(item: Collectable, count: Int): Pocket
  def removeMultipleByType(itemType: CollectableType, count: Int): Pocket
  def contains(item: Collectable): Boolean
  def countByType(itemType: CollectableType): Int
  def isEmpty: Boolean
  def size: Int
\end{lstlisting}

\subsubsection{Generazione e Costruzione del Labirinto}
Parte estensiva del mio lavoro è stata dedicata alla generazione e risoluzione del labirinto,
integrando nel modo più naturale la logica Prolog con Scala.

Il primo aspetto è stato quello di realizzare un costruttore di Maze che potesse assicurare la
corretta generazione del labirinto, soddisfacendo dei requisiti specifici da me imposti.
Infatti, ho reputato che un labirinto potesse essere definito tale solo se:
\begin{itemize}
    \item Il labirinto ha esattamente un ingresso e un'uscita.
    \item L'ingresso e l'uscita sono posti sui bordi del labirinto.
    \item Il labirinto è risolvibile.
\end{itemize}
Per motivi di testing e semplificazione della costruzione di labirinti leggibili, ho deciso di 
prendere esempio dal progetto di Jahrim Gabriele Cesario\cite{pps22chess}, che ha implementato un generatore di
scacchiere. Il mio codice adattato per la DSL (Domain Specific Language) del labirinto è il seguente:
\newpage
\begin{lstlisting}[language=Scala]
object DSL:
    def W(using b: MazeBuilder): MazeBuilder = b + MazeTile.Wall
    def *(using b: MazeBuilder): MazeBuilder = b + MazeTile.Passage
    def E(using b: MazeBuilder): MazeBuilder = b + MazeTile.Entry
    def X(using b: MazeBuilder): MazeBuilder = b + MazeTile.Exit

    extension (builder: MazeBuilder)
      def |(b: MazeBuilder): MazeBuilder = b
\end{lstlisting}
Un esempio di utilizzo della DSL:
\begin{lstlisting}[language=Scala]
val standard: Maze = Maze(10, 11):
    W | W | W | W | E | W | W | W | W | W
    W | * | * | * | * | * | * | * | * | W
    W | * | W | W | W | W | W | W | * | W
    W | * | W | * | * | * | * | W | * | W
    W | * | W | W | * | W | * | W | * | W
    W | * | * | W | * | W | * | * | * | W
    W | W | * | W | * | W | * | W | W | W
    W | * | * | W | * | W | * | W | * | W
    W | * | W | W | W | W | * | W | * | W
    W | * | * | * | * | W | * | * | * | W
    W | W | W | W | X | W | W | W | W | W
\end{lstlisting}

Per la risoluzione del labirinto ho fatto un primo utilizzo di Prolog, sfruttando un semplice 
algoritmo di BFS (Breadth-First Search) per trovare il percorso più breve tra l'ingresso e l'uscita.
Ogni coordinata del labirinto è caricata come fatto, in modo che l'algoritmo possa
riconoscere facilmente in quale direzione può muoversi.
\newpage
\begin{lstlisting}[language=Prolog]
lazy val solution: Option[List[MazePosition]] = computeSolution()

private def computeSolution(): Option[List[MazePosition]] =
  val prologFacts = layout.map((p, t) =>
    t match
      case MazeTile.Passage => s"passage(${p.x}, ${p.y})."
      case MazeTile.Entry => s"entry(${p.x}, ${p.y})."
      case MazeTile.Exit => s"exit(${p.x}, ${p.y})."
      case _ => ""
  ).mkString("\n")
  ...
\end{lstlisting}
Inoltre, il costrutto di Scala "lazy" permette di calcolare la soluzione solamente una volta, evitando
ricalcoli inutili e permettendo di ottenere la soluzione in modo efficiente.

Una volta definite le regole del labirinto ho implementato la logica in Prolog per generarlo 
automaticamente. Non essendoci requisiti specifici sulla difficoltà del labirinto e sulla sua 
struttura, ho deciso di optare per un algoritmo semplice che si adattasse bene allo stile del
Prolog. Per questo, ho implementato il Binary Tree Algorithm, come descritto da Jamis Buck\cite{buck2015mazes},
che richiede una memoria minima e permette un buon livello di casualità e velocità di generazione.
Sebbene non sia il più difficile da risolvere, con grandi dimensioni non diventa banale.
Nello specifico, si tratta di un North-East Tree, e ha questo punto di entrata in Prolog:
\begin{lstlisting}[language=Prolog]
generate_ne_maze(Width, Height, Randoms, RemovedWalls) :-
    Right is Width-1, Bottom is Height-1,
    binary_tree_cells(0, Bottom, Right, Bottom, Randoms, Randoms, RemovedWalls).
\end{lstlisting}
I numeri random per decidere se muoversi a nord o est sono passati come parametro, in quanto
l'algoritmo da me implementato è deterministico e non richiede un generatore di numeri casuali interno, 
non disponibile nella versione di Prolog utilizzata; come valore aggiunto: maggiore è la lunghezza della 
lista di numeri binari, maggiore sarà la difficoltà del labirinto generato.  





